\chapter{Why the Bell Curve Appears}

\begin{center}\itshape
One hundred fair dice sum to $350$ on average. The normal tail beyond $\sqrt{60/7}$\\
is an estimate of the chance that the sum exceeds $400$, and Chapter~5 bounds\\
that normal tail by $0.002$. Hoeffding's inequality bounds the dice by $e^{-2}$.
\end{center}

A single die has mean $7/2$ and variance $35/12$: the second moment is $(1^2+\cdots+6^2)/6=91/6$, and $91/6-(7/2)^2=35/12$. One hundred independent dice have mean $350$ and variance $3500/12$, so the standard deviation is $\sqrt{3500/12}\approx 17.08$. The value $400$ sits $(400-350)/17.08\approx 2.93$ standard deviations above the mean. The number $0.002$ computed below bounds the corresponding normal tail. It does not, by itself, bound the dice. This chapter separates the estimate from a theorem that does.

\begin{figure}[ht]
\centering
\fig{ch06-sums}
\caption{A standardised sum, drawn as a histogram, beside the standard normal density. The central limit theorem says the histogram approaches the curve after centring and scaling. The bars in the picture are schematic; the curve is $\phi$.}
\label{fig:sums}
\end{figure}

\photograph{galton-box}{A wooden Galton board with rows of pins and vertical bins of small balls, taller near the middle.}{A Galton board at the Matemateca of the Institute of Mathematics and Statistics, University of S\~ao Paulo. Photograph by Rodrigo Tetsuo Argenton, 19 December 2016. Credit: Matemateca (IME/USP) / Rodrigo Tetsuo Argenton, Creative Commons Attribution-ShareAlike 4.0, \url{https://commons.wikimedia.org/wiki/File:Galton_box.jpg}. The pins and the bins are the physical experiment. The pile is the balls that were in the machine for this photograph. It is not the schematic histogram above, and it is not a data set for a proof.}

Figure~\ref{fig:sums} is a drawing of the limit. The photograph is one machine, on one day. Galton's quincunx is the reason the limit has a shape you can watch, and the theorem still has to be proved.

\section{Three modes}

\begin{definition}[modes of convergence]\label{def:modes}
Let $X,X_1,X_2,\ldots$ be random variables built from the uniforms of Chapter~5.
\begin{enumerate}[label=(\alph*)]
\item $X_n\to X$ \emph{almost surely} when $P^*(\{X_n\not\to X\})=0$.
\item $X_n\to X$ \emph{in probability} when $P(|X_n-X|>\varepsilon)\to0$ for every $\varepsilon>0$.
\item $X_n\to X$ \emph{in distribution} when $F_n(x)\to F(x)$ at every continuity point $x$ of the distribution function $F$ of $X$.
\end{enumerate}
\end{definition}

\begin{proposition}[comparisons]\label{prop:modes}
Almost-sure convergence implies convergence in probability. Convergence in probability to a constant $c$ implies convergence in distribution to $c$. If $X_n\to c$ in distribution, then $X_n\to c$ in probability.
\end{proposition}
\begin{proof}
On the complement of a set of outer content $0$, $|X_n-X|$ is eventually smaller than $\varepsilon$, so $P(|X_n-X|>\varepsilon)\to0$.

If $X_n\to c$ in probability and $x>c$, then $\{|X_n-c|\le(x-c)/2\}$ is contained in $\{X_n\le x\}$ for the inequality of probabilities up to the event of equality, which does not affect the limit, and that probability tends to $1$. If $x<c$, the same argument gives $F_n(x)\to0$. The constant $c$ has a distribution function that jumps only at $c$.

Conversely, $\{|X_n-c|>\varepsilon\}\subseteq\{X_n\le c-\varepsilon\}\cup\{X_n>c+\varepsilon\}$, so
\[
P(|X_n-c|>\varepsilon)\le F_n(c-\varepsilon)+1-F_n(c+\varepsilon)\to0,
\]
because both endpoints are continuity points of the point mass at $c$. The inclusion avoids an incorrect equality when there is an atom at $c-\varepsilon$.
\end{proof}

\begin{proposition}[test functions]\label{prop:cdf}
If $X_n\to X$ in distribution and $f$ is bounded and uniformly continuous, then $\mathbb E[f(X_n)]\to\mathbb E[f(X)]$. The expectations are Riemann--Stieltjes sums for $f$ against $F_n$ and $F$, taken at continuity points.
\end{proposition}
\begin{proof}
Fix $\varepsilon>0$. Choose $\delta>0$ from uniform continuity so that $|u-v|<\delta$ implies $|f(u)-f(v)|<\varepsilon$. Choose continuity points $a<b$ of $F$ with $F(a)<\varepsilon$ and $1-F(b)<\varepsilon$. For large $n$ the same bounds hold for $F_n$ up to $\varepsilon$. Split the Stieltjes sums into the tails and the compact interval. The tails contribute at most $2\|f\|_\infty\cdot 3\varepsilon$. On $[a,b]$ use a common grid of continuity points spaced by less than $\delta$. On each cell $f$ varies by less than $\varepsilon$, and the $F_n$-mass of the cell differs from the $F$-mass by the discrepancy of $F_n-F$ at the two endpoints, which is small for large $n$.
\end{proof}

\section{A strong law under a fourth moment}

\begin{theorem}[strong law, fourth moment]\label{thm:slln4}
Let $X_i=F^{-1}(U_i)$ be i.i.d.\ copies, in the sense of Propositions~\ref{prop:manyuniforms} and~\ref{prop:inverse}, with $\mathbb E[X_1]=\mu$ and $\mathbb E[X_1^4]<\infty$. Set $S_n=X_1+\cdots+X_n$. Then $S_n/n\to\mu$ almost surely.
\end{theorem}
\begin{proof}
Write $Z_i=X_i-\mu$ and $T_n=Z_1+\cdots+Z_n$. The claim $S_n/n\to\mu$ is the claim $T_n/n\to0$. The proof bounds the outer content of $\{|T_n|\ge n\varepsilon\}$, sums those bounds, and passes to every positive $\varepsilon$ at once.

\emph{The fourth moment is of order $n^2$.} Since $|x-\mu|\le|x|+|\mu|$ and $(a+b)^2\le2(a^2+b^2)$, one has $(a+b)^4\le4(a^2+b^2)^2\le8(a^4+b^4)$. Thus $\mathbb E|Z_1|^4\le8(\mathbb E|X_1|^4+|\mu|^4)<\infty$. Write $\mu_4=\mathbb E[Z_1^4]$ and $\sigma^2=\mathbb E[Z_1^2]$. Expand $T_n^4=\sum_{i,j,k,\ell}Z_i Z_j Z_k Z_\ell$. If any index appears only once, independence and $\mathbb E[Z_i]=0$ make that term's expectation zero. The surviving patterns are a fourth power, and a product of two distinct squares. There are $n$ fourth powers. For an unordered pair $\{i,j\}$ the multinomial coefficient of $Z_i^2 Z_j^2$ is $4!/(2!2!)=6$, and there are $n(n-1)/2$ such pairs, so those terms contribute $6\cdot n(n-1)/2\cdot\sigma^4=3n(n-1)\sigma^4$. Therefore
\[
 \mathbb E[T_n^4]=n\mu_4+3n(n-1)\sigma^4\le Cn^2
\]
for a constant $C$ that does not depend on $n$.

\emph{The deviation event need not be finite-bit.} Each $Z_i$ may depend on infinitely many binary digits, so Borel--Cantelli cannot be quoted for $\{|T_n|\ge n\varepsilon\}$ until an outer-content bound is proved from finite-bit events.

\emph{Clip, then replace the clip by a step function.} Let $Q(u)=F^{-1}(u)-\mu$ on $(0,1)$, so $Z_i=Q(U_i)$, and let $Q_M=\max(-M,\min(M,Q))$. The clipped function is bounded and monotone, so it extends continuously to the compact interval $[0,1]$: if $Q$ is unbounded it is eventually constant near each endpoint, and otherwise the monotone limit is finite. Uniform continuity supplies, for every $\eta>0$, a function $Y_{M,r}$ constant on each dyadic interval of length $2^{-r}$ with $\sup_u|Y_{M,r}(u)-Q_M(u)|<\eta$. The variables $Y_{M,r}(U_i)$ depend on finitely many digits and are independent.

On the event that every $|Z_i|\le M$, the clipped sum $T_{n,M}=\sum_i Q_M(U_i)$ equals $T_n$. If in addition each step is within $\eta$ of the clip, the step sum differs from $T_{n,M}$ by less than $n\eta$. Hence, for $0<\eta<\varepsilon$,
\[
 \{|T_n|\ge n\varepsilon\}
 \subseteq B_M\cup\Bigl\{\Bigl|\sum_{i=1}^n Y_{M,r}(U_i)\Bigr|\ge n(\varepsilon-\eta)\Bigr\},
\]
where $B_M=\{\max_{i\le n}|Z_i|>M\}$. The second event is finite-bit. The first is a union of $n$ events $\{|Q(U_i)|>M\}$. Since $Q$ is increasing, each of those is an initial interval of $[0,1)$ or a final interval, or a union of both. Such an interval can be covered by dyadic cylinders whose total length is arbitrarily close to its own length, so countable subadditivity gives $P^*(B_M)\le n P(|Z_1|>M)$.

\emph{Pass to the limit in the right order.} Markov's inequality on the finite-bit event, followed by the comparison just made, yields
\[
 P^*(|T_n|\ge n\varepsilon)
 \le nP(|Z_1|>M)
 +\frac{\mathbb E\bigl[\bigl(\sum_i Y_{M,r}(U_i)\bigr)^4\bigr]}{n^4(\varepsilon-\eta)^4}.
\]
Fix $M$ and $\eta$. As $r\to\infty$ the step sum converges uniformly to $T_{n,M}$, so the fourth moments converge: choose $r$ large enough that the numerator is at most $\mathbb E[T_{n,M}^4]+n^2$. The elementary bounds $|a+b|^4\le8(|a|^4+|b|^4)$ and $|\sum a_i|^4\le n^3\sum|a_i|^4$, applied to $T_{n,M}=T_n-(T_n-T_{n,M})$ and to the coordinates of the difference, give
\[
 \mathbb E[T_{n,M}^4]
 \le8\mathbb E[T_n^4]+8n^4\mathbb E\bigl[|Z_1|^4\mathbf 1_{\{|Z_1|>M\}}\bigr],
\]
because $|Q-Q_M|\le|Q|$ on $\{|Q|>M\}$ and the difference vanishes elsewhere. Also $P(|Z_1|>M)\le\mathbb E[Z_1^4]/M^4$. Substitute $\mathbb E[T_n^4]\le Cn^2$. The resulting upper bound holds for every $M>0$ and every $\eta\in(0,\varepsilon)$. Send $M\to\infty$ first: both $n\mathbb E[Z_1^4]/M^4$ and the clipped fourth-moment tail tend to $0$. Then send $\eta\downarrow0$. What remains is
\[
 P^*(|T_n|\ge n\varepsilon)\le\frac{8C+1}{n^2\varepsilon^4}.
\]
The extra $1$ is the room left for the dyadic error $n^2$ in the numerator.

\emph{Sum and let $\varepsilon$ tend to zero.} The right-hand side is summable in $n$. The outcomes that violate $|T_n|\ge n\varepsilon$ for infinitely many $n$ sit in $\bigcup_{n\ge N}\{|T_n|\ge n\varepsilon\}$ for every $N$, so their outer content is at most the tail of the series, which tends to $0$. Thus the outer content is $0$ for each fixed $\varepsilon>0$. The countable union over $\varepsilon=1/m$ still has outer content $0$, by Proposition~\ref{prop:subadd}. Outside that union, $T_n/n\to0$, and therefore $S_n/n\to\mu$.
\end{proof}

\section{The empirical distribution}

Let $F_n(x)=n^{-1}\#\{i\le n:X_i\le x\}$.

\begin{theorem}[Glivenko--Cantelli]\label{thm:gc}
For i.i.d.\ copies constructed by the inverse-distribution method, if their common distribution function $F$ is continuous, then $\sup_x|F_n(x)-F(x)|\to0$ almost surely.
\end{theorem}
\begin{proof}
Fix an integer $m\ge2$. Continuity of $F$, the limits $0$ and $1$ at $\pm\infty$, and the intermediate-value theorem supply points $x_1<\cdots<x_{m-1}$ with $F(x_j)=j/m$. At each of these finitely many points, $F_n(x_j)$ is the average of the bounded indicators $\mathbf 1_{\{X_i\le x_j\}}$. In the inverse-distribution model each indicator is a bounded nondecreasing function of one uniform coordinate. The dyadic approximation after Proposition~\ref{prop:inverse} therefore places those averages under the bounded strong law, Theorem~\ref{thm:sllnbound}. There is one null set outside which
\[
 |F_n(x_j)-F(x_j)|\to0\qquad\text{for every }j=1,\ldots,m-1
\]
simultaneously, because a finite union of null sets is null.

On that set, fix $x$ and compare it with the grid. If $x_j\le x<x_{j+1}$, monotonicity of both functions gives $F_n(x_j)\le F_n(x)\le F_n(x_{j+1})$ and $F(x_j)\le F(x)\le F(x_{j+1})$. Hence
\begin{align*}
F_n(x)-F(x)
&\le F_n(x_{j+1})-F(x_j)
= \bigl(F_n(x_{j+1})-F(x_{j+1})\bigr)+\bigl(F(x_{j+1})-F(x_j)\bigr)\\
&\le |F_n(x_{j+1})-F(x_{j+1})|+\frac1m,
\end{align*}
and likewise $F(x)-F_n(x)\le |F_n(x_j)-F(x_j)|+1/m$. If $x<x_1$, then $0\le F(x)\le1/m$ and $0\le F_n(x)\le F_n(x_1)$, so both $F_n(x)-F(x)$ and $F(x)-F_n(x)$ are at most $|F_n(x_1)-F(x_1)|+1/m$. If $x\ge x_{m-1}$, then $F_n(x)-F(x)\le 1-F(x_{m-1})=1/m$, while
\[
 F(x)-F_n(x)\le 1-F_n(x_{m-1})=\frac1m+\bigl(F(x_{m-1})-F_n(x_{m-1})\bigr).
\]
Therefore
\[
 \sup_x|F_n(x)-F(x)|
 \le\frac1m+\max_{1\le j\le m-1}|F_n(x_j)-F(x_j)|.
\]
The maximum tends to $0$, so the limsup of the supremum is at most $1/m$. This is one null set for this fixed $m$. The union of these null sets over $m=1,2,\ldots$ is still null. Outside it, the limsup is at most $1/m$ for every $m$, hence the supremum tends to $0$.
\end{proof}

\section{Characteristic functions}

Write $\varphi_X(t)=\mathbb E[\cos(tX)]+i\,\mathbb E[\sin(tX)]$.

\begin{lemma}[differentiation under the integral]\label{lem:diffint}
If $X$ has a continuous density $f$ and $\int|x|f<\infty$, then $\varphi_X$ is differentiable for every real $t$, with $\varphi'_X(t)=i\mathbb E[Xe^{itX}]$ and $\varphi'_X(0)=i\mathbb E[X]$.
\end{lemma}
\begin{proof}
The mean value theorem gives $|\cos((t+h)x)-\cos(tx)|\le|hx|$. The difference quotients converge uniformly on each bounded $x$-interval by elementary calculus, and the bound $|x|f(x)$ is integrable and continuous. Theorem~\ref{thm:dct} therefore passes the limits inside the integral, separately for sine and cosine. The sine term is the same. At $t=0$ one is left with $\int x f(x)\,\dd x$.
\end{proof}

\begin{proposition}[three transforms]\label{prop:cfexamples}
A constant $c$ has transform $e^{itc}$. A fair sign $\pm1$ has transform $\cos t$. A standard normal has transform $e^{-t^2/2}$.
\end{proposition}
\begin{proof}
The first two are finite sums. For the normal, the sine integral vanishes by oddness. Differentiate $\int\cos(tx)\phi(x)\,\dd x$ by Lemma~\ref{lem:diffint}. Integrate by parts, using $\phi'(x)=-x\phi(x)$, to obtain $\varphi'(t)=-t\varphi(t)$. The solution with $\varphi(0)=1$ is $e^{-t^2/2}$.
\end{proof}

\begin{theorem}[convergence in distribution passes to characteristic functions]\label{thm:cfconv}
If $X_n\to X$ in distribution, then $\varphi_{X_n}(t)\to\varphi_X(t)$ for every fixed $t$.
\end{theorem}
\begin{proof}
The maps $x\mapsto\cos(tx)$ and $x\mapsto\sin(tx)$ are bounded and uniformly continuous. Proposition~\ref{prop:cdf} applies.
\end{proof}

The converse implication is L\'evy's continuity theorem, proved in Volume~11. It is not a step here.

\section{Lindeberg's theorem}

\begin{lemma}[sums of normals]\label{lem:normalsum}
The sum of independent $N(0,\sigma^2)$ and $N(0,\tau^2)$ variables is $N(0,\sigma^2+\tau^2)$, for $\sigma,\tau\ge0$, interpreting variance zero as a constant.
\end{lemma}
\begin{proof}
If either variance is zero, the claim follows by adding the constant zero. Otherwise complete the square in the joint exponent. With $v=\sigma^2+\tau^2$,
\[
\tau^2(x-y)^2+\sigma^2 y^2
=v\Bigl(y-\frac{\tau^2 x}{v}\Bigr)^2+\frac{\sigma^2\tau^2}{v}x^2.
\]
Dividing by $2\sigma^2\tau^2$ separates a Gaussian integral in $y$, of total mass one after the constant in the two densities is included, from the factor $\exp(-x^2/(2v))$. The prefactors produce the constant of the $N(0,v)$ density.
\end{proof}

The general swap below hides its speed. Copies of one variable with a third moment make the speed visible. Figure~\ref{fig:lindeberg} is the bookkeeping: each row replaces one more summand by a normal of the same variance.

\begin{figure}[ht]
\centering
\fig{ch06-lindeberg}
\caption{Six summands, replaced one at a time. The top row is the original sum. Each lower row has turned one more summand into a normal of the same variance. The last row is a sum of normals. The picture is the indexing of the proof, not a sample of values.}
\label{fig:lindeberg}
\end{figure}

\begin{examplex}[one swap]
Let $Y_i$ be i.i.d., centred, with variance $\sigma^2>0$ and $\rho=\mathbb E|Y_1|^3<\infty$. Set $X_{n,k}=Y_k/(\sigma\sqrt n)$ and $S_n=\sum_{k=1}^n X_{n,k}$. Let $\xi_{n,k}$ be independent normals of variance $1/n$, independent of the $Y_i$. Their total variance is $1$, so Lemma~\ref{lem:normalsum} says their sum is standard normal. Let $W_k$ be the sum of $X_{n,1},\ldots,X_{n,k-1}$ and $\xi_{n,k+1},\ldots,\xi_{n,n}$. Then $W_k$ is independent of both $X_{n,k}$ and $\xi_{n,k}$. If $|f'''|\le M$, Taylor's formula with Lagrange remainder gives a remainder $R$ with $|R(x)|\le M|x|^3/6$. Write the expansion at $W_k+X_{n,k}$ and at $W_k+\xi_{n,k}$. The shared value $f(W_k)$ cancels in the difference. The linear terms have mean zero and the quadratic terms match, because the two summands share a mean and a variance and are independent of $W_k$. Each remainder is at most $M|x|^3/6$ at its own increment, so the triangle inequality bounds the difference of expectations by the sum of the two expected remainder bounds. One swap therefore changes $\mathbb E[f]$ by at most
\[
 \frac M6\bigl(\mathbb E|X_{n,k}|^3+\mathbb E|\xi_{n,k}|^3\bigr).
\]
\end{examplex}

\begin{examplex}[the third-moment constant]
The same notation gives $\mathbb E|X_{n,k}|^3=\rho/(\sigma^3 n^{3/2})$, since each $X_{n,k}$ is $Y_k$ divided by $\sigma\sqrt n$. For a standard normal $Z$, start from the density $\phi(x)=(2\pi)^{-1/2}e^{-x^2/2}$, not from the un-normalised integral. The substitution $u=x^2/2$ gives $\int_0^\infty x^3 e^{-x^2/2}\,\dd x=2$, so
\[
 \mathbb E|Z|^3=2\int_0^\infty x^3\phi(x)\,\dd x=\frac{4}{\sqrt{2\pi}}=2\sqrt{\frac{2}{\pi}}.
\]
The two closed forms agree because both are positive and their squares equal $8/\pi$. Each comparison normal is $Z/\sqrt n$, so $\mathbb E|\xi_{n,k}|^3=n^{-3/2}\mathbb E|Z|^3$. Summing the $n$ swaps multiplies by $n/n^{3/2}=n^{-1/2}$. If $f$ has a continuous third derivative bounded by $M$,
\[
 \bigl|\mathbb E[f(S_n)]-\mathbb E[f(Z)]\bigr|
 \le\frac{M}{6\sqrt n}\Biggl(\frac{\rho}{\sigma^3}+2\sqrt{\frac{2}{\pi}}\Biggr).
\]
Replacing every $Y_i$ by $cY_i$ for $c>0$ multiplies $\rho$ by $c^3$ and $\sigma$ by $c$, so $\rho/\sigma^3$ is unchanged, and $S_n$ itself is unchanged. A factor $\sigma^3$ on the Gaussian term would break that invariance. That factor belongs to $Y_i$, whose third moment has not yet been divided by $(\sigma\sqrt n)^3$.
\end{examplex}

Let $H$ be the number of heads in one hundred fair tosses, so $\mathbb EH=50$ and the standard deviation is $5$. Then $P(H\ge 60)=0.02844\ldots$. The normal tail beyond two standard deviations is $1-\Phi(2)=0.02275\ldots$. The continuity correction, at $59.5$, is the tail beyond $1.9$ standard deviations, $0.02872\ldots$. Hoeffding's bound on the same event is $e^{-2}=0.1353\ldots$. The normal number is an approximation. Hoeffding's number is a proved upper bound. The rate above controls smooth test functions of the centred sum $(H-50)/5$. It does not by itself replace either of those two sentences.

\begin{theorem}[Lindeberg]\label{thm:lindeberg}
For each $n$ let $X_{n,1},\ldots,X_{n,n}$ be independent, centred, and satisfy $\sum_k\mathbb E[X_{n,k}^2]=1$. Suppose that for every $\varepsilon>0$,
\[
\sum_{k=1}^n\mathbb E\bigl[X_{n,k}^2\mathbf 1_{|X_{n,k}|>\varepsilon}\bigr]\to0.
\]
Let $S_n=\sum_k X_{n,k}$ and let $Z$ be standard normal. Then $\mathbb E[f(S_n)]\to\mathbb E[f(Z)]$ whenever $f$, $f'$, $f''$ and $f'''$ exist, are continuous, and are bounded.
\end{theorem}
\begin{proof}
Let $\sigma_{n,k}^2=\mathbb E[X_{n,k}^2]$, and let $\xi_{n,k}$ be independent normals of those variances, independent of the array. Lemma~\ref{lem:normalsum}, applied repeatedly, says that $\sum_k\xi_{n,k}$ is standard normal. Swap one coordinate at a time. Let $A_{n,k}$ be the sum of $X_{n,1},\ldots,X_{n,k-1}$ and $\xi_{n,k+1},\ldots,\xi_{n,n}$. Taylor's formula with Lagrange remainder gives
\[
f(A+x)=f(A)+f'(A)x+\tfrac12 f''(A)x^2+R(x),
\]
where $|R(x)|\le M|x|^3/6$ if $|f'''|\le M$. If $|f''|\le B$, the expansion stopped one order earlier also yields $|R(x)|\le B x^2$. Thus $|R(x)|\le\min(M|x|^3/6,\,Bx^2)$.

The linear terms have expectation $0$ and the quadratic terms have the same expectation, because $X_{n,k}$ and $\xi_{n,k}$ share a mean and a variance and are independent of $A_{n,k}$. The difference of the expectations of $f$ is at most the sum of the expected remainders. On $|x|\le\varepsilon$ the remainder is at most $(M\varepsilon/6)x^2$. On $|x|>\varepsilon$ it is at most $B x^2$. Summing on $k$, the pieces that carry a factor $\varepsilon$ contribute at most $M\varepsilon/3$, since both arrays have total variance $1$. The pieces that carry the indicator of $\{|X_{n,k}|\ge\varepsilon\}$ tend to $0$ by hypothesis.

For the Gaussian pieces, $\max_k\sigma_{n,k}^2\to0$. Indeed
\[
\sigma_{n,k}^2=\mathbb E[X^2\mathbf 1_{|X|\le\varepsilon}]+\mathbb E[X^2\mathbf 1_{|X|>\varepsilon}]\le\varepsilon^2+\mathbb E[X^2\mathbf 1_{|X|>\varepsilon}],
\]
so the largest variance is at most $\varepsilon^2$ plus the Lindeberg sum. A normal of variance $\sigma^2$ has fourth moment $\mathbb E[\xi^4]=3\sigma^4$, by two integrations by parts using $\phi'(x)=-x\phi(x)$. Therefore
\[
\sum_k\mathbb E[\xi_{n,k}^2\mathbf 1_{|\xi_{n,k}|>\varepsilon}]
\le\frac1{\varepsilon^2}\sum_k\mathbb E[\xi_{n,k}^4]
=\frac3{\varepsilon^2}\sum_k\sigma_{n,k}^4
\le\frac3{\varepsilon^2}(\max_k\sigma_{n,k}^2)\sum_k\sigma_{n,k}^2\to0.
\]
This corrects the tempting but false inequality that the truncated Gaussian second moment is bounded by variance times the Gaussian tail probability.

The limsup of the difference of expectations is therefore at most a constant times $\varepsilon$. Let $\varepsilon\downarrow0$.
\end{proof}

\begin{corollary}[classical central limit theorem]\label{cor:clt}
Let $Y_i$ be i.i.d.\ with mean $\mu$ and variance $\sigma^2\in(0,\infty)$, and put $S_n=Y_1+\cdots+Y_n$. Then $(S_n-n\mu)/(\sigma\sqrt n)$ converges in distribution to $N(0,1)$.
\end{corollary}
\begin{proof}
Apply Theorem~\ref{thm:lindeberg} to $X_{n,k}=(Y_k-\mu)/(\sigma\sqrt n)$. The sum of the variances is $1$, and
\[
\sum_k\mathbb E[X_{n,k}^2\mathbf 1_{|X_{n,k}|>\varepsilon}]
=\sigma^{-2}\mathbb E[(Y_1-\mu)^2\mathbf 1_{|Y_1-\mu|>\varepsilon\sigma\sqrt n}]\to0,
\]
because $(Y_1-\mu)^2$ is integrable. Thus $\mathbb E[f(S_n)]\to\mathbb E[f(Z)]$ for every admissible test function $f$. To pass from test functions to distribution functions, fix a continuity point $x$ of $\Phi$ and a width $\delta>0$. Use explicit smooth cutoffs rather than Gaussian convolution, which cannot stay between sharp indicator functions. Define $q(u)=1$ for $u\le0$, $q(u)=0$ for $u\ge1$, and
\[
q(u)=1-35u^4+84u^5-70u^6+20u^7\qquad(0<u<1).
\]
Its first three derivatives vanish at both endpoints, so $q$ is $C^3$ with bounded derivatives. For $\delta>0$, the functions $f_-(y)=q((y-x+\delta)/\delta)$ and $f_+(y)=q((y-x)/\delta)$ satisfy
\[
f_-(y)\le\mathbf1_{\{y\le x\}}\le f_+(y).
\]
Lindeberg applies to both. Their normal expectations lie between $\Phi(x-\delta)$ and $\Phi(x+\delta)$ in the required directions, giving
\[
\Phi(x-\delta)\le\liminf_nF_n(x)\le\limsup_nF_n(x)\le\Phi(x+\delta).
\]
Let $\delta\downarrow0$ and use continuity of $\Phi$.
\end{proof}

One hundred dice fall under this corollary. The centred sum, in units of its standard deviation, exceeds
\[
 \frac{400-350}{\sqrt{3500/12}}=\sqrt{\frac{60}7},
\]
because $50^2/(3500/12)=2500\cdot12/3500=60/7$. Chapter~5's Gaussian tail bounds the probability by $\phi(x)/x$ at $x=\sqrt{60/7}$. The exponent is $x^2/2=30/7=4+2/7$. From $e^{-1}<1/2.7$, as in the coin bound of Chapter~4,
\[
 e^{-4}<(1/2.7)^4=1/53.1441<0.01882.
\]
For $0<y<1$, the Taylor remainder gives $e^{-y}<1-y+y^2/2$. At $y=2/7$,
\[
 1-\frac27+\frac{2}{49}=\frac{37}{49}<0.756.
\]
Thus $e^{-30/7}<0.01882\cdot0.756<0.01423$. Also $2\pi>6.25$, so $1/\sqrt{2\pi}<0.4$, and $x^2=60/7>8.41=2.9^2$, so $x>2.9$. Therefore
\[
 \frac{\phi(x)}{x}<\frac{0.01423\cdot0.4}{2.9}<0.002.
\]
The inequality $\phi(x)/x<0.002$ bounds the standard normal tail. It is an estimate of the dice only in the sense that Corollary~\ref{cor:clt} says the centred and scaled sum converges in distribution to that normal law. Convergence in distribution does not supply a numerical error at $n=100$, so $0.002$ is not a proved upper bound on $P(S\ge 400)$.

Hoeffding's inequality is a bound on the dice. Each summand lies in $[1,6]$, an interval of length $5$. Theorem~\ref{thm:hoeffding}, with $t=50$ and $\sum_i(b_i-a_i)^2=100\cdot25=2500$, gives
\[
 P(S\ge 400)=P(S-\mathbb E S\ge 50)\le\exp\Bigl(-\frac{2\cdot 50^2}{2500}\Bigr)=e^{-2}.
\]
Chapter~4 records $e^{-2}=0.1353\ldots$. The normal estimate is smaller than $0.002$. The proved bound is $e^{-2}$. They answer different questions, and the smaller number is not a substitute for the larger one.

\section{The local theorem}

\begin{theorem}[local de Moivre--Laplace]\label{thm:local}
Let $S_n$ be binomial with parameters $n$ and $1/2$. If integers $k_n$ satisfy $(2k_n-n)/\sqrt n\to y$, then
\[
 P(S_n=k_n)\sim\sqrt{\frac{2}{\pi n}}\exp\Bigl(-\frac{(2k_n-n)^2}{2n}\Bigr).
\]
In particular $P(S_{2m}=m)\sim 1/\sqrt{\pi m}$.
\end{theorem}
\begin{proof}
Stirling's formula from Volume~1, Chapter~8, gives $\ln m!=m\ln m-m+\tfrac12\ln(2\pi m)+o(1)$ as $m\to\infty$. Set $k=k_n$ and $p=k/n$, and assume $p$ stays in a neighbourhood of $1/2$, which it does because $(2k-n)/\sqrt n$ stays bounded. Then
\begin{align*}
\ln P(S_n=k)
&= \ln n!-\ln k!-\ln(n-k)!-n\ln 2\\
&= -k\ln(2k/n)-(n-k)\ln(2(n-k)/n)+\tfrac12\ln\Bigl(\frac{n}{2\pi k(n-k)}\Bigr)+o(1).
\end{align*}
The logarithm of the square-root prefactor tends to the logarithm of $\sqrt{2/(\pi n)}$, because $k(n-k)/n^2\to1/4$. For the power terms, set $\delta=(2k-n)/(2n)$, so $p=1/2+\delta$, $2p=1+2\delta$, $2(1-p)=1-2\delta$, and $\delta=O(n^{-1/2})$. Then
\[
 p\ln(2p)+(1-p)\ln(2(1-p))=p\ln(1+2\delta)+(1-p)\ln(1-2\delta).
\]
The expansion $\ln(1+u)=u-u^2/2+u^3/3+O(u^4)$, with $u=\pm2\delta$, gives
\begin{align*}
\ln(1+2\delta)&=2\delta-2\delta^2+\tfrac83\delta^3+O(\delta^4),\\
\ln(1-2\delta)&=-2\delta-2\delta^2-\tfrac83\delta^3+O(\delta^4).
\end{align*}
Multiply by $p=1/2+\delta$ and by $1-p=1/2-\delta$ and add. The terms in $\delta$ cancel, the terms in $\delta^3$ cancel, and the $\delta^2$ terms contribute $2\delta^2$. The error is $O(\delta^4)$. Multiply by $n$:
\[
 n\cdot 2\delta^2=\frac{(2k-n)^2}{2n},
\]
while $n\delta^4=O(1/n)\to0$. In the formula for $\ln P(S_n=k)$ this combination is subtracted, so the power contribution is $-(2k-n)^2/(2n)+o(1)$. Exponentiating gives the claim. The central case $n=2m$, $k=m$ is $y=0$, and $\sqrt{2/(\pi\cdot 2m)}=1/\sqrt{\pi m}$.
\end{proof}

\section{Gaussian vectors, and cycles}

\begin{proposition}[cycles as Bernoulli summands]\label{prop:cyclebern}
The number $C_n$ of cycles of a uniform random permutation of $\{1,\ldots,n\}$ has the same law as $\sum_{k=1}^n B_k$, where the $B_k$ are independent and $P(B_k=1)=1/k$.
\end{proposition}
\begin{proof}
Insert the symbol $n$ into a uniform random permutation of $\{1,\ldots,n-1\}$. There are $n$ equally likely slots in the cycle notation: after any of the existing symbols, or in a new cycle. Exactly one slot opens a new cycle, and the choice is independent of the smaller permutation. Thus $C_n=C_{n-1}+B_n$. Induct from $C_1=1$.
\end{proof}

\begin{theorem}[Gaussian vectors]\label{thm:gaussvec}
Let $Z$ be a vector of $d$ independent standard normals, and let $A$ be a real $d\times d$ matrix. Then $X=AZ$ has covariance $\Sigma=AA^{\mathsf T}$. If $\Sigma$ is invertible, $X$ has density
\[
 (2\pi)^{-d/2}(\det\Sigma)^{-1/2}\exp\bigl(-\tfrac12 x^{\mathsf T}\Sigma^{-1}x\bigr).
\]
The coordinates of a centred Gaussian vector are jointly independent if and only if every pair of distinct coordinates has covariance zero.
\end{theorem}
\begin{proof}
$\mathbb E[X_i X_j]$ expands into $\Sigma_{ij}$ because $\mathbb E[Z_a Z_b]=\mathbf 1_{a=b}$. The spectral theorem in Volume~1, Chapter~7, factorises the positive definite matrix $\Sigma$ as $QDQ^{\mathsf T}$ with $D$ diagonal and positive. Then $A_0=QD^{1/2}$ satisfies $A_0A_0^{\mathsf T}=\Sigma$, and $A_0^{-1}X=(A_0^{-1}A)Z$. The square matrix $O=A_0^{-1}A$ obeys $OO^{\mathsf T}=A_0^{-1}\Sigma(A_0^{-1})^{\mathsf T}=I$, so it is orthogonal. An orthogonal change of coordinates preserves the density $(2\pi)^{-d/2}\exp(-|z|^2/2)$: the integrand depends only on $|z|^2$, and the absolute value of the determinant, which Volume~1 identifies with the volume scaling of the unit cube, equals $1$. Returning to $x$ multiplies by $|\det A_0|^{-1}=(\det\Sigma)^{-1/2}$ and replaces $|z|^2$ by $x^{\mathsf T}\Sigma^{-1}x$.

For the independence assertion, suppose $\Sigma$ is diagonal, allowing zero diagonal entries. The row vectors $a_i$ of $A$ are then pairwise orthogonal, since $a_i\cdot a_j=\Sigma_{ij}=0$ for $i\ne j$. Normalize each nonzero row and extend these orthonormal vectors to an orthonormal basis of $\R^d$. The associated orthogonal matrix $O$ preserves the product standard-normal density by the same change-of-coordinates argument above. Each nonconstant $X_i=a_i\cdot Z$ is a scalar multiple of a \emph{distinct} coordinate of $OZ$, so all $X_i$ are jointly independent; a zero row gives a constant coordinate and causes no exception. Conversely joint independence and centring give $\mathbb E[X_iX_j]=\mathbb E[X_i]\mathbb E[X_j]=0$ for every $i\ne j$.
\end{proof}

The covariance $\min(s,t)$ of Chapter~12 is the case of this theorem along a finite set of times.

\section{Exercises}

\begin{exercise}[$\star\star$]\label{ex:slutsky}
Suppose $X_n\to X$ in distribution and $Y_n\to c$ in probability. Show that $X_n+Y_n\to X+c$ and $X_n Y_n\to cX$ in distribution.
\end{exercise}

\begin{exercise}[$\star$]\label{ex:confidence}
Let $Y_i$ be i.i.d.\ with mean $\mu$, variance $\sigma^2>0$ and a finite fourth moment. Let $\hat\sigma_n^2=n^{-1}\sum_i(Y_i-\bar Y_n)^2$. Show that
\[
 P\bigl(|\bar Y_n-\mu|\le 1.96\,\hat\sigma_n/\sqrt n\bigr)\to 2\Phi(1.96)-1.
\]
\end{exercise}

\begin{exercise}[$\star\star$]\label{ex:delta}
In the setting of Exercise~\ref{ex:confidence}, let $g$ be continuously differentiable near $\mu$ with $g'(\mu)\ne0$. Show that $\sqrt n(g(\bar Y_n)-g(\mu))$ converges in distribution to $N(0,g'(\mu)^2\sigma^2)$.
\end{exercise}

\begin{exercise}[$\star$]\label{ex:cauchyclt}
Let $Y_i$ have the Cauchy density of Exercise~\ref{ex:cauchy}, independently. Show that $S_n/\sqrt n$ does not converge in distribution, and say which hypothesis of Corollary~\ref{cor:clt} fails.
\end{exercise}

\begin{exercise}[$\star\star$]\label{ex:poissonhalf}
Let $N_n$ be Poisson with mean $n$. Show that $P(N_n\le n)\to1/2$.
\end{exercise}

\begin{exercise}[$\star\star$]\label{ex:returns}
Let $S_{2m}=2B-2m$ with $B$ binomial$(2m,1/2)$. Show that $P(S_{2m}=0)\sim 1/\sqrt{\pi m}$.
\end{exercise}

\begin{exercise}[$\star\star$]\label{ex:thirdmoment}
For a standard normal $Z$, prove $\mathbb E|Z|^3=2\sqrt{2/\pi}$ from the substitution in the third-moment example. For a fair coin, the centred summand is $\pm 1$, so $\rho=\sigma=1$. Write the numerical value of the smooth-test bound at $n=100$ when $M=1$.
\end{exercise}

\begin{exercise}[$\star\star\star$]\label{ex:cyclesclt}
Let $C_n$ be the number of cycles in a uniform random permutation of $n$ elements. Show that $(C_n-\ln n)/\sqrt{\ln n}$ converges in distribution to $N(0,1)$.
\end{exercise}

\section{Hints}

\begin{hint}{ex:slutsky}
Compare $f(X_n+Y_n)$ with $f(X_n+c)$ where $f$ is bounded and uniformly continuous.
\end{hint}
\begin{hint}{ex:confidence}
The fourth moment makes $Y_1^2$ eligible for the weak law. Then use Exercise~\ref{ex:slutsky}.
\end{hint}
\begin{hint}{ex:delta}
$g(\bar Y_n)-g(\mu)=(g'(\mu)+\varepsilon_n)(\bar Y_n-\mu)$ with $\varepsilon_n\to0$ in probability.
\end{hint}
\begin{hint}{ex:cauchyclt}
Prove that the characteristic function is $e^{-|t|}$ from the identity $1/(1+x^2)=\int_0^\infty e^{-u}\cos(ux)\,\dd u$, then let the damping $\varepsilon\downarrow0$ in the switched integral. The product of $n$ transforms, evaluated at $t/\sqrt n$, is $\exp(-\sqrt n\,|t|)$.
\end{hint}
\begin{hint}{ex:poissonhalf}
Corollary~\ref{cor:clt} at the continuity point $0$ of $\Phi$.
\end{hint}
\begin{hint}{ex:returns}
Theorem~\ref{thm:local} at $y=0$, or Stirling's formula on $\binom{2m}{m}$.
\end{hint}
\begin{hint}{ex:thirdmoment}
The integral $\int_0^\infty x^3 e^{-x^2/2}\,\dd x$ equals $2$. Multiply by the constant in $\phi$ and double it for the two tails. The bound in the example is then a single arithmetic line.
\end{hint}
\begin{hint}{ex:cyclesclt}
Proposition~\ref{prop:cyclebern} and Theorem~\ref{thm:lindeberg}. The variance is $H_n-\sum k^{-2}\sim\ln n$.
\end{hint}

\section{Solutions}

\begin{solution}{ex:slutsky}
\idea A uniformly continuous function barely moves when its argument moves a little, and the probability of a large move tends to zero.
\writeup Let $f$ be bounded and uniformly continuous. Given $\varepsilon$, choose $\delta$ so that $|u-v|<\delta$ implies $|f(u)-f(v)|<\varepsilon$. Then the expectations of $f(X_n+Y_n)$ and $f(X_n+c)$ differ by at most $\varepsilon+2\|f\|_\infty P(|Y_n-c|\ge\delta)$. The second term tends to $0$, and $\mathbb E[f(X_n+c)]\to\mathbb E[f(X+c)]$ by Proposition~\ref{prop:cdf}. For the product, $X_n(Y_n-c)\to0$ in probability: the family $X_n$ is tight because it converges in distribution, so $P(|X_n|>A)$ is small uniformly in large $n$, and on $|X_n|\le A$ the product is at most $A|Y_n-c|$. Add this null term to $cX_n$ by the additive case.
\end{solution}

\begin{solution}{ex:confidence}
\idea Slutsky's lemma replaces the unknown standard deviation by the sample one.
\writeup Corollary~\ref{cor:clt} gives the limit $N(0,1)$ for the studentised mean with $\sigma$ in the denominator. Finite fourth moment implies that $Y_1^2$ has finite variance, so the weak law yields $n^{-1}\sum Y_i^2\to\mathbb E[Y_1^2]$ in probability. Also $\bar Y_n\to\mu$ in probability, so $\hat\sigma_n^2\to\sigma^2$ and $\hat\sigma_n\to\sigma$ in probability. Exercise~\ref{ex:slutsky} replaces $\sigma$ by $\hat\sigma_n$. The limit distribution function is continuous, so the probability of $[-1.96,1.96]$ tends to $2\Phi(1.96)-1$.
\end{solution}

\begin{solution}{ex:delta}
\idea Differentiability writes the increment as the derivative times the increment of the mean, plus a smaller term.
\writeup Continuity of $g'$ at $\mu$ and $\bar Y_n\to\mu$ in probability give
\[
 g(\bar Y_n)-g(\mu)=\bigl(g'(\mu)+\varepsilon_n\bigr)(\bar Y_n-\mu)
\]
with $\varepsilon_n\to0$ in probability. Multiply by $\sqrt n$. The main term converges in distribution to $N(0,g'(\mu)^2\sigma^2)$. The error is $\varepsilon_n$ times a tight sequence, so it tends to $0$ in probability. Exercise~\ref{ex:slutsky} adds the two pieces.
\end{solution}

\begin{solution}{ex:cauchyclt}
\idea The characteristic function of a Cauchy variable is $e^{-|t|}$. At scale $1/\sqrt n$ it collapses to $0$ for every $t\ne0$, which no genuine limit law can do.
\writeup Integration by parts gives $\int_0^\infty e^{-u}\cos(ux)\,\dd u=1/(1+x^2)$. For $\varepsilon>0$ the integral
\[
 I_\varepsilon(t)=\int_{-\infty}^{\infty}\frac{e^{-\varepsilon|x|}\cos(tx)}{1+x^2}\,\dd x
\]
equals $\int_0^\infty e^{-u}\bigl(\int_{-\infty}^{\infty}e^{-\varepsilon|x|}\cos(tx)\cos(ux)\,\dd x\bigr)\dd u$, by Proposition~\ref{prop:fubini} on compact ranges of $u$ and the integrable dominant $e^{-u}e^{-\varepsilon|x|}$. The inner integral equals
\[
 \frac{\varepsilon}{\varepsilon^2+(t+u)^2}+\frac{\varepsilon}{\varepsilon^2+(t-u)^2}.
\]
Substitute $u=t+\varepsilon v$ in the second term, for $t>0$. It becomes $\int e^{-(t+\varepsilon v)}/(1+v^2)\,\dd v$, with limits expanding to $\R$ as $\varepsilon\downarrow0$. The integrand is dominated by an integrable function of $v$ for small $\varepsilon$, on the region $\varepsilon|v|\le t/2$, and the far tails are small. Theorem~\ref{thm:dct} passes the limit and yields $\pi e^{-t}$. The first term tends to $0$. Dividing by $\pi$ shows that the characteristic function is $e^{-|t|}$.

Independence multiplies characteristic functions, so $S_n/\sqrt n$ has transform $\exp(-\sqrt n\,|t|)$. If the laws converged, Theorem~\ref{thm:cfconv} would force the transforms to converge to a characteristic function, which equals $1$ at $t=0$ and is continuous there by the same dominated estimate as Lemma~\ref{lem:diffint} with dominant $1$. The limit $\exp(-\sqrt n\,|t|)$ is $0$ for each $t\ne0$, which cannot be continuous at $0$ unless it is identically $0$. So there is no convergence in distribution. The missing hypothesis in Corollary~\ref{cor:clt} is a finite positive variance. Exercise~\ref{ex:cauchy} shows $\mathbb E|Y_1|=\infty$. On $\{|Y_1|\ge1\}$ one has $Y_1^2\ge|Y_1|$, so the second moment is infinite as well.
\end{solution}

\begin{solution}{ex:poissonhalf}
\idea The Poisson law of mean $n$ has variance $n$, so the central limit theorem applies at the point $0$.
\writeup The Poisson-addition result of Chapter~2 shows that Poisson$(n)$ has the same law as the sum of $n$ independent Poisson$(1)$ variables. Each has mean and variance $1$, so Corollary~\ref{cor:clt} applies to the fixed i.i.d.\ law and gives $(N_n-n)/\sqrt n\to N(0,1)$ in distribution. Then $P(N_n\le n)=P((N_n-n)/\sqrt n\le 0)\to\Phi(0)=1/2$.
\end{solution}

\begin{solution}{ex:returns}
\idea The middle binomial coefficient is Stirling's formula at its most symmetric point.
\writeup Theorem~\ref{thm:local} with $n=2m$ and $k=m$ gives $P(S_{2m}=0)\sim 1/\sqrt{\pi m}$. The same conclusion is Stirling's formula alone:
\[
 \binom{2m}{m}\sim\frac{4^m}{\sqrt{\pi m}}.
\]
\end{solution}

\begin{solution}{ex:thirdmoment}
\idea One Gaussian integral fixes the constant. The coin is the case $\sigma=\rho=1$.
\writeup Let $\phi(x)=(2\pi)^{-1/2}e^{-x^2/2}$. Then $\int_0^\infty x^3 e^{-x^2/2}\,\dd x=2$, by $u=x^2/2$. So
\[
 \mathbb E|Z|^3=2\int_0^\infty x^3\phi(x)\,\dd x=2\cdot\frac{2}{\sqrt{2\pi}}=2\sqrt{\frac{2}{\pi}}.
\]
For the fair coin, $\rho/\sigma^3=1$. The Gaussian term is $2\sqrt{2/\pi}$, with no extra factor of $\sigma^3$, because the comparison normals have variance $1/n$. With $M=1$ and $n=100$ the bound is
\[
 \frac{1}{60}\Biggl(1+2\sqrt{\frac{2}{\pi}}\Biggr)=\frac{1+1.5957\ldots}{60}=0.04326\ldots.
\]
The same number is obtained if each centred toss is replaced by $c$ times itself: both $\rho$ and $\sigma^3$ scale by $c^3$.
\end{solution}

\begin{solution}{ex:cyclesclt}
\idea Independent Bernoulli summands of mean $1/k$ have variance $\ln n$, and none of them is large compared with that standard deviation.
\writeup Proposition~\ref{prop:cyclebern} gives $C_n\stackrel{d}=\sum_{k=1}^n B_k$. Then $\mathbb E[C_n]=H_n$ and $\operatorname{Var}C_n=H_n-\sum_{k\le n}k^{-2}\sim\ln n$, by Theorem~\ref{thm:coupon} and the convergence of $\sum k^{-2}$. Each centred summand is at most $1$ in absolute value, so the normalised summand is at most $1/\sqrt{\operatorname{Var}C_n}$. Once $\sqrt{\operatorname{Var}C_n}>1/\varepsilon$, every indicator $\{|X_{n,k}|>\varepsilon\}$ is empty and the Lindeberg sum is $0$. Theorem~\ref{thm:lindeberg} and the passage in Corollary~\ref{cor:clt} give convergence after centring at $H_n$ and scaling by the standard deviation. Exercise~\ref{ex:slutsky} replaces both $H_n$ and the variance by $\ln n$, since their ratio tends to $1$ and their difference is bounded.
\end{solution}

\section*{Sources and further reading}

\begin{enumerate}
\item A.~de Moivre, \emph{Approximatio ad summam terminorum binomii}, privately printed in 1733 and reprinted in \emph{The Doctrine of Chances}, 1738. The local normal approximation for a fair binomial.
\item J.~W.~Lindeberg, Eine neue Herleitung des Exponentialgesetzes in der Wahrscheinlichkeitsrechnung, \emph{Math.\ Z.} \textbf{15} (1922), 211--225. Theorem~\ref{thm:lindeberg}.
\item V.~Glivenko, Sulla determinazione empirica delle leggi di probabilit\`a, \emph{Giorn.\ Ist.\ Ital.\ Attuari} \textbf{4} (1933), 92--99. Theorem~\ref{thm:gc}, with Cantelli's note in the same volume, 421--424.
\item Volume~1, Chapter~8, for Stirling's formula.
\end{enumerate}

